% Chapter 1 content — included by main.tex; compile chapter1.tex for preview.
\section{Introduction}
\label{sec:introduction}

In recent years, non-geostationary satellite orbit (NGSO) systems, especially low Earth orbit (LEO) constellations, have become an important component of broadband satellite communications, non-terrestrial networks (NTN), and future 6G architectures~\cite{giordani2021ntn,lin2021leosat6g,zhujiang2022istn,alhraishawi2022ngsoSurvey}. Unlike geostationary satellite orbit (GSO) systems, which are located above the Earth's equator and appear fixed with respect to the ground, LEO satellites move rapidly along their orbits and can provide lower latency and more flexible coverage. However, when NGSO downlink transmissions operate in the same or adjacent frequency bands as existing GSO systems, they may cause co-channel interference to GSO ground stations. In particular, when a LEO satellite passes close to the line-of-sight direction between a GSO ground station and its target GSO satellite, a near-inline geometry may occur, leading to a higher interference risk, as illustrated in Fig.~\ref{fig:ch1_coexistence_scenario}.

\begin{figure}[t]
  \centering
  \includegraphics[width=0.95\columnwidth]{{ch1_NGSO-GSO co-existence scenrio}}
  \caption{NGSO--GSO coexistence scenario and near-inline interference geometry between a LEO satellite, a GSO ground station, and the target GSO satellite.}
  \label{fig:ch1_coexistence_scenario}
\end{figure}

To protect existing GSO systems from interference generated by NGSO systems, the International Telecommunication Union (ITU) specifies interference limits for NGSO systems in Article~22 of the Radio Regulations~\cite{itur_rr2020}. Among these limits, equivalent power flux-density (EPFD) is used to evaluate the aggregate interference generated by NGSO transmissions toward GSO systems. Since EPFD depends on satellite geometry, transmit power, antenna gain, propagation loss, and the off-axis angle observed at the GSO ground station, its value changes over time as LEO satellites move along their orbits. Therefore, an NGSO system must ensure that the aggregate EPFD received at the GSO ground station remains below the applicable EPFD limit.

To satisfy the EPFD constraint, NGSO systems may adopt different interference mitigation strategies, such as transmit power control, satellite tilt or pitch adjustment, beam shutdown, and beam steering~\cite{hills2022beamsteering,jalali2023joint,li2021progressive}. These approaches can reduce the interference observed at the GSO ground station by lowering the transmitted power toward the GSO direction, changing the satellite pointing geometry, or avoiding high-risk beam transmissions. However, such mitigation actions may also reduce the available service resources for LEO users. For example, reducing beam power can decrease the achievable downlink rate, while shutting down high-risk beams can directly interrupt the service of users originally covered by those beams. Therefore, NGSO--GSO coexistence should not only be evaluated from the perspective of EPFD compliance, but also from the perspective of service continuity and user satisfaction.

In a multi-beam LEO system, the service impact of EPFD mitigation becomes more critical during near-inline events. When a LEO satellite passes close to the GSO ground station and the GSO line-of-sight direction, its beams may contribute a large portion of the aggregate EPFD. To satisfy the EPFD limit, the system may need to shut down some high-risk beams on the dominant interfering satellite. Although beam shutdown is effective in reducing the EPFD contribution toward the GSO ground station, it also interrupts the service of users originally served by those closed beams. Therefore, the users under closed beams become the main service demand that must be recovered during the critical period.

This service recovery problem creates an opportunity for beam reassociation. LEO satellite motion and beam footprints are predictable, so the coverage relationship among neighboring satellites can be estimated in advance. In particular, same-orbit neighboring satellites may provide sufficient coverage overlap with the dominant interfering satellite during a critical period. Therefore, users affected by closed-beam shutdown may still be covered by neighboring LEO beams and can potentially be reassociated to these beams for service recovery. This observation motivates the use of neighboring satellites as helper satellites for EPFD-aware service recovery through beam reassociation.

Based on this observation, this work proposes EPFD-Aware Beam Reassociation for Service Recovery (EABR), a beam reassociation framework for multi-beam LEO satellite systems under NGSO--GSO coexistence. The main idea is to exploit the predictable coverage overlap among neighboring LEO satellites. When high-risk beams are shut down to satisfy the EPFD limit, users originally served by those beams may still be covered by nearby helper beams. The proposed EABR framework therefore uses beam reassociation to transfer affected users from closed beams to feasible helper recovery beams, so that their service can be continuously supported during the critical period. In this way, EABR aims to improve user satisfaction while maintaining EPFD compliance.

The main contributions of this work are summarized as follows.

\begin{itemize}[leftmargin=*,labelindent=1.5em,nosep]
  \item We propose EPFD-Aware Beam Reassociation for Service Recovery (EABR), a beam reassociation framework for EPFD-constrained multi-beam LEO systems. EABR exploits predictable coverage overlap among neighboring LEO satellites to continuously support users affected by critical-beam shutdown.

  \item We design the core components of the proposed framework, including a time-slot beam role record, Safe-Beam Reassociation for Helper Users (SBR), released power allocation for helper recovery beams, and Helper-Beam Reassociation for Closed-Beam Users (HBR). The time-slot beam role record captures critical satellites, helper satellites, and overlap-based candidate beam relation graphs. SBR uses critical safe beams to reassociate selected helper users and release helper satellite transmit power. The released power is then allocated to helper recovery beams, while HBR reassociates closed-beam users to feasible helper recovery beams.

  \item We evaluate the proposed framework using an STK--MATLAB simulation framework~\cite{agi2023stk,mathworks2023matlab}. The results show that the proposed framework satisfies the EPFD constraint while improving average user satisfaction, enhancing service recovery capability, and reducing the proportion of low-satisfaction users compared with baseline methods.
\end{itemize}
